% ========================================================
% BÀI TẬP TOÁN 5 - KẾT NỐI TRI THỨC VỚI CUỘC SỐNG
% Bài 74: Ôn tập một số yếu tố thống kê và xác suất
% Định dạng: BTTeX Part III - Trắc nghiệm điền số ngắn
% ========================================================

% Câu 1
\begin{ex}
  Bảng thống kê số học sinh giỏi của 4 lớp khối 5: 5A: 18 bạn, 5B: 20 bạn, 5C: 22 bạn, 5D: 20 bạn. Lớp có nhiều học sinh giỏi nhất có bao nhiêu bạn?
  \shortans{22}
  \loigiai{
    Lớp 5C có 22 bạn (nhiều nhất).
  }
\end{ex}

% Câu 2
\begin{ex}
  Ở bảng trên, tổng số học sinh giỏi của khối 5 là bao nhiêu bạn?
  \shortans{80}
  \loigiai{
    $18 + 20 + 22 + 20 = 80$ bạn.
  }
\end{ex}

% Câu 3
\begin{ex}
  Ở bảng trên, trung bình mỗi lớp có bao nhiêu học sinh giỏi?
  \shortans{20}
  \loigiai{
    $80 : 4 = 20$ bạn.
  }
\end{ex}

% Câu 4
\begin{ex}
  Biểu đồ quạt tròn biểu diễn thể loại sách mượn tại thư viện: Khoa học $35\%$, Văn học $40\%$, Lịch sử $25\%$. Tỉ số phần trăm sách Khoa học và Lịch sử cộng lại là bao nhiêu phần trăm?
  \shortans{60}
  \loigiai{
    $35\% + 25\% = 60\%$.
  }
\end{ex}

% Câu 5
\begin{ex}
  Nếu có 400 lượt mượn sách ở biểu đồ trên thì số lượt mượn sách Văn học là bao nhiêu lượt?
  \shortans{160}
  \loigiai{
    $400 \times 40\% = 160$ lượt.
  }
\end{ex}

% Câu 6
\begin{ex}
  Ở bài trên, số lượt mượn sách Lịch sử là bao nhiêu lượt?
  \shortans{100}
  \loigiai{
    $400 \times 25\% = 100$ lượt.
  }
\end{ex}

% Câu 7
\begin{ex}
  Tung một đồng xu 40 lần, có 18 lần xuất hiện mặt ngửa. Số lần xuất hiện mặt sấp là bao nhiêu lần?
  \shortans{22}
  \loigiai{
    $40 - 18 = 22$ lần.
  }
\end{ex}

% Câu 8
\begin{ex}
  Ở thực nghiệm trên, tỉ số phần trăm số lần xuất hiện mặt ngửa là bao nhiêu phần trăm?
  \shortans{45}
  \loigiai{
    $(18 : 40) \times 100\% = 45\%$.
  }
\end{ex}

% Câu 9
\begin{ex}
  Gieo một con xúc xắc 50 lần, xuất hiện mặt 6 chấm 12 lần. Tỉ số phần trăm số lần xuất hiện mặt 6 chấm là bao nhiêu phần trăm?
  \shortans{24}
  \loigiai{
    $(12 : 50) \times 100\% = 24\%$.
  }
\end{ex}

% Câu 10
\begin{ex}
  Trong hộp có 5 viên bi xanh, 3 viên bi đỏ, 2 viên bi vàng. Tổng số viên bi trong hộp là bao nhiêu viên?
  \shortans{10}
  \loigiai{
    $5 + 3 + 2 = 10$ viên.
  }
\end{ex}

% Câu 11
\begin{ex}
  Ở hộp bi trên, tỉ số phần trăm số bi xanh so với tổng số bi là bao nhiêu phần trăm?
  \shortans{50}
  \loigiai{
    $(5 : 10) \times 100\% = 50\%$.
  }
\end{ex}

% Câu 12
\begin{ex}
  Ở hộp bi trên, tỉ số phần trăm số bi đỏ so với tổng số bi là bao nhiêu phần trăm?
  \shortans{30}
  \loigiai{
    $(3 : 10) \times 100\% = 30\%$.
  }
\end{ex}

% Câu 13
\begin{ex}
  Biểu đồ quạt tròn biểu diễn tỉ lệ học sinh tham gia các câu lạc bộ: Vẽ $25\%$, Nhạc $30\%$, Thể thao $45\%$. Câu lạc bộ nào có đông học sinh nhất chiếm bao nhiêu phần trăm?
  \shortans{45}
  \loigiai{
    Thể thao chiếm nhiều nhất với 45%.
  }
\end{ex}

% Câu 14
\begin{ex}
  Nếu có 200 học sinh tham gia ở câu trên thì câu lạc bộ Vẽ có bao nhiêu bạn?
  \shortans{50}
  \loigiai{
    $200 \times 25\% = 50$ bạn.
  }
\end{ex}

% Câu 15
\begin{ex}
  Thực hiện quay vòng tròn may mắn 100 lần, có 35 lần kim chỉ ô may mắn. Tỉ số phần trăm ô may mắn xuất hiện là bao nhiêu phần trăm?
  \shortans{35}
  \loigiai{
    $35\%$.
  }
\end{ex}

% Câu 16
\begin{ex}
  Bảng thống kê lượng mưa 3 tháng hè (mm): Tháng 6: 180, Tháng 7: 240, Tháng 8: 300. Trung bình mỗi tháng lượng mưa là bao nhiêu mm?
  \shortans{240}
  \loigiai{
    $(180 + 240 + 300) : 3 = 720 : 3 = 240\text{ mm}$.
  }
\end{ex}

% Câu 17
\begin{ex}
  Ở bài trên, tháng mưa nhiều nhất hơn tháng mưa ít nhất bao nhiêu mm?
  \shortans{120}
  \loigiai{
    $300 - 180 = 120\text{ mm}$.
  }
\end{ex}

% Câu 18
\begin{ex}
  Một tổ học sinh gồm 6 bạn có số điểm kiểm tra: 8, 9, 10, 8, 9, 10. Điểm trung bình của tổ là bao nhiêu điểm?
  \shortans{9}
  \loigiai{
    $(8 + 9 + 10 + 8 + 9 + 10) : 6 = 54 : 6 = 9$ điểm.
  }
\end{ex}

% Câu 19
\begin{ex}
  Khảo sát ý kiến của 80 người tiêu dùng, có 64 người đánh giá tốt sản phẩm. Tỉ số phần trăm người đánh giá tốt là bao nhiêu phần trăm?
  \shortans{80}
  \loigiai{
    $(64 : 80) \times 100\% = 80\%$.
  }
\end{ex}

% Câu 20
\begin{ex}
  Ở khảo sát trên, số người không đánh giá tốt là bao nhiêu người?
  \shortans{16}
  \loigiai{
    $80 - 64 = 16$ người.
  }
\end{ex}

% Câu 21
\begin{ex}
  Một hình tròn trên biểu đồ quạt chia làm 3 phần: Phần A chiếm $40\%$, Phần B chiếm $35\%$. Phần C chiếm bao nhiêu phần trăm?
  \shortans{25}
  \loigiai{
    $100\% - 40\% - 35\% = 25\%$.
  }
\end{ex}

% Câu 22
\begin{ex}
  Gieo xúc xắc 50 lần, có 22 lần ra mặt chấm lẻ. Số lần ra mặt chấm chẵn là bao nhiêu lần?
  \shortans{28}
  \loigiai{
    $50 - 22 = 28$ lần.
  }
\end{ex}

% Câu 23
\begin{ex}
  Ở bài trên, tỉ số phần trăm ra mặt chấm chẵn là bao nhiêu phần trăm?
  \shortans{56}
  \loigiai{
    $(28 : 50) \times 100\% = 56\%$.
  }
\end{ex}

% Câu 24
\begin{ex}
  Tung 1 đồng xu 10 lần, có 7 lần ngửa. Tỉ số số lần ngửa so với tổng số lần tung nhân với 10 là bao nhiêu?
  \shortans{7}
  \loigiai{
    $\frac{7}{10} \times 10 = 7$.
  }
\end{ex}

% Câu 25
\begin{ex}
  Số học sinh của 3 lớp 5A, 5B, 5C lần lượt là 32, 34, 36 bạn. Trung bình mỗi lớp có bao nhiêu bạn?
  \shortans{34}
  \loigiai{
    $(32 + 34 + 36) : 3 = 102 : 3 = 34$ bạn.
  }
\end{ex}
